% ======================================================================
%  PRED001 — The Actualization Theory: Prediction Registry
%  A Reconstruction of Physics from Difference, Actualization and Spectrum
%
%  Standalone publication-grade document for Zenodo.
%
%  Status: COMPLETE — canonical-derivation registry
%  Method: assembly only — no new physics, no new derivations, no new fits
%  Inputs: canonical monograph V2.0 (Ch1-14) and the claim-classification
%          registry (ATQG_ClaimClassificationRegistry.md), the immutable
%          prediction registry (ATQG_Predictions.md, QG193) and the
%          prediction-outcome dashboard (ATQG_PredictionOutcomes.md, QG202).
%
%  Build: pdflatex ActualizationTheory_PredictionRegistry.tex (twice)
%  Rules: accepted canonical derivations only; no new physics; no new fits;
%         no speculative predictions; no retrospective reinterpretation.
% ======================================================================

\documentclass[11pt,a4paper]{article}

% ---- Fonts and encoding ------------------------------------------------
\usepackage[T1]{fontenc}
\usepackage[utf8]{inputenc}
\usepackage{lmodern}

% ---- Mathematics --------------------------------------------------------
\usepackage{amsmath,amssymb}

% ---- Links and colors ----------------------------------------------------
\usepackage[hidelinks]{hyperref}
\usepackage{xcolor}

% ---- Page geometry --------------------------------------------------------
\usepackage[a4paper,margin=2.5cm]{geometry}
\usepackage{setspace}
\onehalfspacing

% ---- Header/footer --------------------------------------------------------
\usepackage{fancyhdr}
\pagestyle{fancy}
\fancyhf{}
\fancyhead[L]{\small The Actualization Theory --- Prediction Registry}
\fancyhead[R]{\small Version 1.0}
\fancyfoot[C]{\thepage}
\renewcommand{\headrulewidth}{0.4pt}

% ---- Table of contents depth ---------------------------------------------
\setcounter{tocdepth}{2}

% ---- Prediction record environment ---------------------------------------
% \predstart{ID}{STATUS} opens a record; \pfield{Label}{Content} fills a
% labelled line. Records are separated by rules and may span pages.
\newcommand{\predstart}[2]{%
  \par\medskip\noindent\hrule\par\smallskip
  \noindent{\large\bfseries #1}\hfill{\bfseries Status: #2}\par\smallskip
}
\newcommand{\pfield}[2]{%
  \noindent\textbf{#1:}~#2\par\smallskip
}

% ---- Canonical D96 constants (declared once, used throughout) ------------
\newcommand{\Sigmam}{95}
\newcommand{\SigmaSqrtm}{64.08}
\newcommand{\SigmaMtwo}{229}
\newcommand{\occMom}{1900.25}
\newcommand{\Dspan}{6.40}
\newcommand{\occZero}{4}
\newcommand{\occTwo}{87}

% ======================================================================
\begin{document}

% ---- Title ----------------------------------------------------------------
\thispagestyle{empty}
\begin{center}
  \vspace*{1.5cm}
  {\Huge\bfseries The Actualization Theory}\\[0.6em]
  {\LARGE A Reconstruction of Physics from\\ Difference, Actualization and Spectrum}\\[1.2em]
  {\LARGE\bfseries Prediction Registry}\\[2.2em]
  {\Large Fabrice Wieser}\\[0.5em]
  {\small Independent Researcher}\\[0.3em]
  {\small Senior Software Engineer}\\[0.3em]
  {\small \url{https://github.com/MagusDraconis/AT}}\\[2.0em]
  {\large The Actualization Theory (AT) --- Prediction Registry v1.0}\\[0.5em]
  {\large Companion document to The Actualization Theory (AT) v2.0 (Canonical)}\\[0.5em]
  {\large DOI: 10.5281/zenodo.20681734}\\[0.5em]
  {\large \today}
  \vfill
\end{center}
\clearpage

% ---- Abstract ---------------------------------------------------------------
\begin{abstract}
\noindent
This document is the official \emph{Prediction Registry} of The Actualization Theory
(AT): a standalone, publication-grade registry of the quantitative and qualitative
predictions derived from the canonical theory, compiled before future experimental
confirmation. Every entry is derived from the accepted canonical derivations of the
monograph (Difference $\to$ Actualization $\to$ Inevitable Spectrum $\to$ Physics), carries
its derivation source, its dependency chain, its numerical value (where applicable), its
uncertainty, its validation status, and its falsification criterion. The registry uses no
new physics, no new fits, and no retrospective reinterpretation; it records the canonical
content as published. Statuses distinguish \textsc{Confirmed} (pre-registered target
verified), \textsc{Consistent} (value matches current observation), \textsc{Pending}
(awaiting validation), and \textsc{Falsified} (contradicted by observation). No registered
prediction is currently falsified.
\end{abstract}

\clearpage

\tableofcontents
\clearpage

% ======================================================================
%  Part I — Introduction
% ======================================================================
\part{Introduction}

\section{Purpose of the Prediction Registry}
\label{sec:purpose}

The Prediction Registry exists to fix, in advance of experiment, the falsifiable content of
The Actualization Theory. Its purposes are:

\begin{itemize}
  \item to record every quantitative and qualitative prediction derived from the canonical
        theory, with its derivation source and dependency chain;
  \item to fix the numerical values and uncertainty ranges before future confirmation, so
        that confirmation or contradiction is genuine independent evidence and not post-hoc
        fitting;
  \item to document, per prediction, the experimental test, the current evidence, and the
        falsification criterion;
  \item to serve as the immutable, version-controlled reference for external validation
        (the theory's pre-registration program, QG190--QG193, and its outcome dashboard,
        QG202);
  \item to be suitable for independent Zenodo publication and future version-controlled
        updates.
\end{itemize}

The registry is \emph{assembly only}: it contains no new physics, no new derivations, no new
fits, and no speculation. Every entry is drawn from the accepted canonical results of the
monograph V2.0 and the claim-classification registry.

\section{Validation Philosophy}
\label{sec:philosophy}

The Actualization Theory distinguishes \emph{derivation} from \emph{validation}
(monograph Ch.~14, Theorem~c14:thm:derived-vs-validated): a quantity may be derived without
being validated, and its derivation is complete before its validation is possible. This
registry records derived predictions and their validation status on that basis.

The statuses used throughout are:

\begin{itemize}
  \item \textsc{Confirmed}: a pre-registered target verified by observation (the prediction
        was frozen before the data).
  \item \textsc{Consistent}: the derived value matches current observation within the
        stated deviation (including calibrated correspondences and documented fits). These
        are retrospective consistency records, not pre-registered confirmations.
  \item \textsc{Pending}: awaiting validation --- the current experimental frontier.
  \item \textsc{Falsified}: contradicted by observation. As of this version, no registered
        prediction is falsified.
\end{itemize}

The classification of each prediction's \emph{derivational} status --- theorem, necessity,
correspondence, calibration, hosted, fit --- is preserved from the claim-classification
registry and disclosed in each record. In particular, entries classified as
\emph{calibration} carry an anchor ($v$, $m_e$, or SI conversion) and entries classified as
\emph{fit} are labelled as fitted readouts, not derivations.

\section{Pre-Registration Principles}
\label{sec:pre-registration}

The theory maintains an immutable prediction registry (QG193) containing three
pre-registered predictions with frozen values and explicit falsification conditions:
P1 (a 106 GeV resonance), P2 (the $0\nu\beta\beta$ effective Majorana mass), and P3 (the
sector-ladder spectrum). The pre-registration principles are:

\begin{itemize}
  \item \textbf{Frozen values}: each pre-registered prediction's value is fixed before
        future data; no future phase may modify it (registry lock, QG193).
  \item \textbf{Forbidden inputs}: pre-registered predictions assert that no experimental
        limit, detector sensitivity, fitted mass, or future measurement entered the
        derivation (forbidden-input guards, QG190--QG192).
  \item \textbf{State advance only}: the outcome state may advance forward
        (PENDING $\to$ SUPPORTED $\to$ CONFIRMED, or PENDING $\to$ DISFAVORED $\to$
        FALSIFIED); frozen values never change.
  \item \textbf{Temporal independence}: confirmation of a pre-registered prediction is
        genuine independent evidence because the value was frozen before the observation.
\end{itemize}

The three pre-registered predictions are registered in full in Section~\ref{sec:frontier}.
All other entries in this registry are derived values with documented consistency or pending
status.

% ======================================================================
%  Part II — Standard Model Predictions
% ======================================================================
\part{Standard Model Predictions}

\section{Gauge Sector}
\label{sec:gauge}

\predstart{AT-P001}{CONSISTENT}
\pfield{Title}{Gauge structure $1+3+8$ (dimensional correspondence).}
\pfield{Derived From}{Difference $\to$ Actualization $\to$ Spectrum (D96) $\to$ degree-12
connectivity of $C_{96}(\pm1,\ldots,\pm6)$.}
\pfield{Prediction}{The gauge sector of the standard model, $U(1)\times SU(2)\times SU(3)$,
corresponds dimensionally to the D96 sector structure: one $U(1)$, three $SU(2)$, and eight
$SU(3)$ generators, $1+3+8 = 12$, the degree of the $C_{96}$ generator ring. The gauge
groups and their Lie algebras are \emph{hosted}, not derived (claim classification:
dimensional correspondence).}
\pfield{Experimental Test}{Count of gauge generators in the observed standard model.}
\pfield{Current Evidence}{The standard model carries exactly 12 gauge generators
(1 photon + 3 weak + 8 gluons). The correspondence is dimensional; the groups are hosted.}
\pfield{Future Validation Path}{No experimental change expected; the claim is a dimensional
correspondence, not a prediction of new group content.}
\pfield{Falsification Criterion}{Observation of a gauge sector not matching the $1+3+8$
dimensional count (not currently conceivable at low energy).}

\predstart{AT-P002}{CONSISTENT}
\pfield{Title}{Gauge couplings as spectral ratios.}
\pfield{Derived From}{Difference $\to$ Actualization $\to$ Spectrum (D96) $\to$ spectral
constants $\Sigma m$, $\Sigma\sqrt{m}$, $\#d$; couplings
$1/\alpha_{\mathrm{em}} = \Sigma m + \#d = 137$ (fitted readout),
$\alpha_{\mathrm{weak}} = 3/\Sigma m$, $\alpha_{\mathrm{strong}} = 8/\Sigma\sqrt{m}$,
$\sin^2\theta_W = 0.2316$.}
\pfield{Prediction}{The weak and strong couplings correspond to D96 spectral ratios
($3/95$, $8/64.08$); the fine-structure inverse is the fitted readout $137 = 95+42$.
None of the couplings carries a defined renormalization scale (claim classification:
$1/\alpha_{\mathrm{em}}$ = fit; $\alpha_{\mathrm{weak}}$, $\alpha_{\mathrm{strong}}$ =
correspondence).}
\pfield{Experimental Test}{Measurement of the gauge couplings and the Weinberg angle.}
\pfield{Current Evidence}{$\sin^2\theta_W = 0.2316$ matches the measured value; the weak and
strong ratios match the observed couplings at the natural scale; $1/\alpha_{\mathrm{em}} =
137$ matches the low-energy fine-structure constant (documented fit).}
\pfield{Future Validation Path}{Precision running of $\alpha_s$ to the $Z$ pole:
spectral $\alpha_s = 8/\Sigma\sqrt{m} = 0.1248$ vs.\ PDG $\alpha_s(M_Z) = 0.1184$
(dev.\ $5.4\%$); improvement of the measured low-energy $\alpha_{\mathrm{em}}$.}
\pfield{Falsification Criterion}{A measured gauge coupling inconsistent with the spectral
ratios at the natural scale beyond the stated correspondence accuracy.}

\predstart{AT-P003}{CONSISTENT}
\pfield{Title}{Weak boson masses from the weak scale.}
\pfield{Derived From}{Difference $\to$ Actualization $\to$ Spectrum (D96) $\to$ weak scale
$v = (\Sigma m + \#d)\ln(\mathrm{span}) = 254$ GeV.}
\pfield{Prediction}{$M_W = 80.1$ GeV and $M_Z = 91.4$ GeV from the Weinberg projection, with
$\rho = M_W^2/(M_Z^2\cos^2\theta_W) = 1$ exactly at tree level.}
\pfield{Experimental Test}{Measurement of the $W$ and $Z$ masses (LEP/ATLAS/CMS).}
\pfield{Current Evidence}{Measured $M_W \simeq 80.4$ GeV, $M_Z = 91.19$ GeV; the
reconstruction agrees within the calibration accuracy (claim classification: calibration via
the anchor $v$).}
\pfield{Future Validation Path}{Future $W$-mass refinements at HL-LHC; $M_Z$ is precisely
known.}
\pfield{Falsification Criterion}{A measured $M_W$/$M_Z$ ratio inconsistent with $\rho = 1$
beyond the calibration accuracy.}

\predstart{AT-P004}{CONSISTENT}
\pfield{Title}{Precision electroweak observables.}
\pfield{Derived From}{Difference $\to$ Actualization $\to$ Spectrum (D96) $\to$ spectral
constants; $\sin^2\theta_{\mathrm{eff}} = \#g/(2\Sigma m) = 0.23158$;
$\Gamma_Z = M_H\cos\theta_W/\#g = 2.4953$; $\Gamma_W = \sigma_{\mathrm{occ}}^2/
(\mathrm{occMom}\cdot\lambda_2) = 2.0852$; $\Gamma_H = \lambda_2/\Sigma m = 4.07$ MeV;
$R_b = \mathrm{span}\cdot g_2\cdot\sin^4\theta_W = 0.2163$;
$A_{FB}^b = (\lambda_H/\lambda_2)^2 = 0.0992$; $A_{FB}^\ell = M_H/(M_W M_Z) = 0.0171$.}
\pfield{Prediction}{The precision electroweak observables (effective Weinberg angle, $Z$
and $W$ widths, Higgs width, $R_b$, forward-backward asymmetries) are read as spectral
ratios (claim classification: correspondence).}
\pfield{Experimental Test}{Measurements of the $Z$-pole observables, $R_b$, and the
asymmetries at LEP/SLC and LHC.}
\pfield{Current Evidence}{$\sin^2\theta_{\mathrm{eff}} = 0.23158$ matches the measured
$0.23153$ within $0.03\%$; the widths and $R_b$ match within the stated deviations
(0.004\%--0.05\%).}
\pfield{Future Validation Path}{HL-LHC precision runs; any improvement of the $R_b$ and
asymmetry measurements.}
\pfield{Falsification Criterion}{A measured precision observable deviating from its spectral
ratio beyond the stated accuracy.}

\section{CKM and PMNS}
\label{sec:mixing}

\predstart{AT-P005}{CONSISTENT}
\pfield{Title}{CKM matrix as octave-transition read.}
\pfield{Derived From}{Difference $\to$ Actualization $\to$ Spectrum (D96) $\to$ octave
transitions: $V_{us} = \#d/(2\Sigma m)$,
$V_{cb} = (\Omega_0/\Omega_2)^{\delta_d}$, $V_{ub} = 2V_{cb}\cdot(\mathrm{occ}_0/
\mathrm{occ}_2)$; $\Omega_0 = 0.924$, $\Omega_2 = 3.385$ are the octave-family centers
(not the Laplacian modes $\omega_k = \sqrt{\lambda_k}$), $\delta_d = 2.449$ the down-sector
effective dimension.}
\pfield{Prediction}{The quark mixing matrix elements are spectral correspondences matched to
the measured CKM matrix with mean deviation $0.58\%$ (claim classification:
correspondence).}
\pfield{Experimental Test}{Measurement of the CKM elements in weak decays.}
\pfield{Current Evidence}{The three reconstructed elements match the measured values with
mean deviation $0.58\%$.}
\pfield{Future Validation Path}{Improvement of $V_{cb}$ and $V_{ub}$ measurements (Belle~II,
LHCb); the ratios are fixed spectral forms.}
\pfield{Falsification Criterion}{A measured CKM element deviating from its spectral ratio
beyond the stated correspondence accuracy.}

\predstart{AT-P006}{CONSISTENT}
\pfield{Title}{CKM CP violation and the Jarlskog invariant.}
\pfield{Derived From}{Difference $\to$ Actualization $\to$ Spectrum (D96) $\to$ chiral
circulation: $\sin\delta = \mathrm{occ}_{\mathrm{top}}/\Sigma m = 0.9158$,
$\delta_{CP} = 66.3^\circ$, $J = 3.139\times10^{-5}$.}
\pfield{Prediction}{The CKM CP phase and the Jarlskog invariant are spectral-ratio reads
(claim classification: correspondence).}
\pfield{Experimental Test}{Measurement of CP-violating observables in the quark sector.}
\pfield{Current Evidence}{$\delta_{CP} \simeq 66^\circ$ is consistent with the global CKM
fit; $J = 3.139\times10^{-5}$ matches the measured $J \simeq 3\times10^{-5}$.}
\pfield{Future Validation Path}{Precision CP measurements at Belle~II and LHCb.}
\pfield{Falsification Criterion}{A measured $J$ inconsistent with the spectral value beyond
the stated accuracy.}

\predstart{AT-P007}{CONSISTENT}
\pfield{Title}{PMNS matrix as T3-only read.}
\pfield{Derived From}{Difference $\to$ Actualization $\to$ Spectrum (D96) $\to$ T3-only
access reads: $\theta_{12} = 33.35^\circ$, $\theta_{23} = 49.72^\circ$,
$\theta_{13} = 8.34^\circ$, $\delta_\nu = 66.4^\circ$, mean deviation $1.5\%$.}
\pfield{Prediction}{The lepton mixing angles and the Dirac CP phase are spectral
correspondences (claim classification: correspondence).}
\pfield{Experimental Test}{Measurement of the PMNS angles in neutrino oscillations.}
\pfield{Current Evidence}{The angles match the measured values with mean deviation $1.5\%$.}
\pfield{Future Validation Path}{Improved $\theta_{23}$ and $\delta_\nu$ determination
(DUNE, JUNO, Hyper-K); the $\delta_\nu = 66.4^\circ$ read is a continuous spectral-ratio
phase, not lattice-restricted.}
\pfield{Falsification Criterion}{A measured PMNS angle deviating from its spectral read
beyond the stated correspondence accuracy.}

\section{Neutrino Sector}
\label{sec:neutrino}

\predstart{AT-P008}{CONSISTENT}
\pfield{Title}{Neutrino mass splittings.}
\pfield{Derived From}{Difference $\to$ Actualization $\to$ Spectrum (D96) $\to$ closed-form
laws: $\Delta m^2_{21} = (1/\Sigma\sqrt{m})^2/(\mathrm{span}/2) = 7.607\times10^{-5}$
eV$^2$ (dev.\ $1.02\%$), $\Delta m^2_{31} = \sin^2\theta_W/\Sigma m = 2.44\times10^{-3}$
eV$^2$ (dev.\ $0.71\%$).}
\pfield{Prediction}{The solar and atmospheric splittings are closed-form D96 ratios,
calibrated to eV$^2$ by an implicit scale (claim classification: correspondence +
calibration).}
\pfield{Experimental Test}{Measurement of the splittings in solar, atmospheric, and reactor
neutrino oscillations.}
\pfield{Current Evidence}{$\Delta m^2_{21} = 7.607\times10^{-5}$ and
$\Delta m^2_{31} = 2.44\times10^{-3}$ eV$^2$ match the measured values within the stated
deviations.}
\pfield{Future Validation Path}{Precision oscillation experiments (JUNO, DUNE, Hyper-K).}
\pfield{Falsification Criterion}{A measured splitting deviating from its closed-form ratio
beyond the stated accuracy.}

\predstart{AT-P009}{PENDING}
\pfield{Title}{Absolute neutrino masses (normal ordering).}
\pfield{Derived From}{Difference $\to$ Actualization $\to$ Spectrum (D96) $\to$ absolute
mass law: $m_1 = 0$ (zero-mode, normal ordering), $m_2 = 1/(\Sigma\sqrt{m}\cdot\sqrt{
\mathrm{span}/2}) = 8.7216$ meV (phys.\ $8.72$), $m_3 = \sqrt{\#g}/(\Sigma m\sqrt{2}) =
49.3728$ meV (phys.\ $49.4$); $\Sigma m_\nu = 0.0581$ eV.}
\pfield{Prediction}{The absolute neutrino masses are closed-form D96 expressions with the
normal ordering $m_1 = 0 < m_2 < m_3$ and $\Sigma m_\nu = 0.0581$ eV $< 0.12$ eV
(claim classification: correspondence + calibration).}
\pfield{Experimental Test}{Direct mass measurement (KATRIN, next-generation tritium
beta-decay) and cosmology (sum of masses from CMB/LSS).}
\pfield{Current Evidence}{$\Sigma m_\nu = 0.0581$ eV is below current limits
($\Sigma m_\nu \lesssim 0.12$ eV from cosmology); the absolute scale is not yet measured;
the mass ordering is not yet determined.}
\pfield{Future Validation Path}{KATRIN sensitivity reach; JUNO mass-ordering determination;
DESI/Euclid sum-of-masses bounds.}
\pfield{Falsification Criterion}{A measured $\Sigma m_\nu$ inconsistent with $0.0581$ eV, or
an inverted ordering established ($m_3 < m_2$).}

\predstart{AT-P010}{PENDING}
\pfield{Title}{Neutrino mass ordering (normal).}
\pfield{Derived From}{Same chain as AT-P009: the zero-mode mass $m_1 = 0$ fixes the normal
ordering.}
\pfield{Prediction}{The neutrino mass ordering is normal ($m_1 = 0 < m_2 < m_3$).}
\pfield{Experimental Test}{Mass-ordering determination via long-baseline oscillations
(DUNE, Hyper-K) and matter effects.}
\pfield{Current Evidence}{Current data mildly prefer normal ordering; not conclusive.}
\pfield{Future Validation Path}{DUNE (first half of the 2030s) and Hyper-K.}
\pfield{Falsification Criterion}{Establishment of inverted ordering with significance.}

\predstart{AT-P011}{CONSISTENT}
\pfield{Title}{Quark mass ratios (absolute masses).}
\pfield{Derived From}{Difference $\to$ Actualization $\to$ Spectrum (D96) $\to$ unified
mass law: $m_u = m_e\cdot\Sigma\sqrt{m}/\sqrt{\Sigma m^2} = 2.164$ (dev.\ $0.18\%$);
$m_d = m_u\cdot(\Sigma\sqrt{m})^2/\mathrm{occMom} = 4.676$ (dev.\ $0.14\%$);
$m_s = m_d\cdot\mathrm{occMom}/\Sigma m = 93.5$; $m_c = m_d\cdot(\Sigma\sqrt{m})^2/
\sqrt{\Sigma m^2} = 1269$; $m_b = m_d\cdot\mathrm{occMom}^2\cdot\Sigma m\cdot\#g/
(\Sigma\sqrt{m})^4 = 4186$; $m_t = m_u\cdot\mathrm{occMom}\cdot\#d = 172704$ MeV.}
\pfield{Prediction}{All six absolute quark masses are reconstructed from the electron anchor
and D96 spectral moments, each within $0.2\%$ (claim classification: calibration via the
anchor $m_e$).}
\pfield{Experimental Test}{Measurement of the quark masses in the $\overline{\mathrm{MS}}$
scheme.}
\pfield{Current Evidence}{All six masses match the PDG $\overline{\mathrm{MS}}$ values
within $0.2\%$.}
\pfield{Future Validation Path}{Lattice QCD refinements of the light-quark masses.}
\pfield{Falsification Criterion}{A measured quark mass deviating from its spectral
construction beyond $0.2\%$.}

\section{Majorana Character}
\label{sec:majorana}

\predstart{AT-P012}{PENDING}
\pfield{Title}{Neutrino Majorana character.}
\pfield{Derived From}{Difference $\to$ Actualization $\to$ Spectrum (D96) $\to$ T3-only
self-conjugate channel (48/95 modes), unique $Q=0$ sector with no conserved charge, real
mass matrix from the reflection automorphism.}
\pfield{Prediction}{The neutrino is Majorana: self-conjugate in its $T_3$-only channel, with
no conserved charge and a real mass matrix (claim classification: correspondence).}
\pfield{Experimental Test}{Observation of neutrinoless double beta decay ($0\nu\beta\beta$).}
\pfield{Current Evidence}{No $0\nu\beta\beta$ observation; current limits
$m_{\beta\beta} \lesssim 0.036$--$0.156$ eV are far above the predicted $2.02$ meV.}
\pfield{Future Validation Path}{nEXO and LEGEND-1000 ton-scale experiments.}
\pfield{Falsification Criterion}{A conclusive Dirac-character determination (e.g.\ a
fundamental conservation law excluding $0\nu\beta\beta$).}

\predstart{AT-P013}{PENDING}
\pfield{Title}{Effective Majorana mass $m_{\beta\beta}$ (pre-registered P2).}
\pfield{Derived From}{Difference $\to$ Actualization $\to$ Spectrum (D96) $\to$ PMNS reads
(AT-P007) + absolute masses (AT-P009) + Majorana character (AT-P012): $m_{\beta\beta} =
\bigl|m_1c_{12}^2c_{13}^2 + m_2s_{12}^2c_{13}^2e^{i\alpha_2} +
m_3s_{13}^2e^{-2i\delta_\nu}\bigr| = 2.02$ meV (computed $2.0222$ meV), with
$\alpha_2 = \alpha_3 = 0$.}
\pfield{Prediction}{$m_{\beta\beta} = 2.02$ meV, dominated by the term
$m_2 s_{12}^2 c_{13}^2 = 2.52$ meV and robust to the CP phase (pre-registered, QG191).}
\pfield{Experimental Test}{Search for $0\nu\beta\beta$ at nEXO / LEGEND-1000 (half-life
sensitivity $\sim10^{28}$ yr).}
\pfield{Current Evidence}{None --- the prediction is below all current limits
($0.036$--$0.156$ eV).}
\pfield{Future Validation Path}{nEXO / LEGEND-1000 ton-scale reach.}
\pfield{Falsification Criterion}{A measured upper limit below $2.02$ meV (significant
exclusion) falsifies P2.}

\section{Electron g-2}
\label{sec:g2}

\predstart{AT-P014}{CONSISTENT}
\pfield{Title}{Electron anomalous magnetic moment.}
\pfield{Derived From}{Difference $\to$ Actualization $\to$ Spectrum (D96) $\to$ Schwinger
term modulated by the lightest-octave fraction:
$a_e = (\alpha/2\pi)\bigl(1 - (\mathrm{occ}_0/\Sigma m)^2\bigr) = 1.159655\times10^{-3}$
(dev.\ $0.0003\%$).}
\pfield{Prediction}{The electron anomaly is the Schwinger term with the negative spectral
correction $1 - (\mathrm{occ}_0/\Sigma m)^2$ (claim classification: correspondence).}
\pfield{Experimental Test}{Precision measurement of $a_e$ (Penning-trap experiments) and
theoretical QED computation.}
\pfield{Current Evidence}{Matches experiment within $0.0003\%$.}
\pfield{Future Validation Path}{Improved QED tenth-order terms and fine-structure
determination.}
\pfield{Falsification Criterion}{A measured $a_e$ inconsistent with the spectral value
beyond the stated accuracy.}

\predstart{AT-P015}{CONSISTENT}
\pfield{Title}{Electron g-2 anomaly-free beyond $\Delta a_e$.}
\pfield{Derived From}{Same chain: $\Delta a_e = (\alpha/2\pi)^3\,\mathrm{span}^{1/4}\,
(\mathrm{occ}_0/\Sigma m)^3 = 1.86\times10^{-13} < 10^{-12}$.}
\pfield{Prediction}{No electron g-2 anomaly: the new-physics contribution is below
$10^{-12}$ (claim classification: correspondence).}
\pfield{Experimental Test}{Comparison of measured $a_e$ with the SM prediction at
$10^{-12}$ precision.}
\pfield{Current Evidence}{No electron g-2 anomaly observed; the value matches SM QED.}
\pfield{Future Validation Path}{Continued improvement of $\alpha$ and the QED series.}
\pfield{Falsification Criterion}{An electron g-2 discrepancy exceeding
$\Delta a_e = 1.86\times10^{-13}$.}

\predstart{AT-P016}{CONSISTENT}
\pfield{Title}{Muon g-2 (spectral D96 mechanism).}
\pfield{Derived From}{Difference $\to$ Actualization $\to$ Spectrum (D96) $\to$ same
Schwinger-plus-spectral-fraction mechanism with the dense-bulk sign:
$a_\mu = (\alpha/2\pi)(1 + \lambda_2/\Sigma m) = 1.16644\times10^{-3}$ (dev.\ $0.045\%$);
$\Delta a_\mu = (\alpha/2\pi)^3\,\mathrm{span}^{1/4} = 2.494\times10^{-9}$ (dev.\ $0.24\%$).}
\pfield{Prediction}{The muon anomaly is the same D96 mechanism as the electron, with the
positive dense-bulk correction $+\lambda_2/\Sigma m$ (opposite spectral end)
(claim classification: correspondence).}
\pfield{Experimental Test}{Measurement of $a_\mu$ at Fermilab (Muon g-2) and J-PARC.}
\pfield{Current Evidence}{The central value matches within $0.045\%$; the anomaly magnitude
is consistent with the measured g-2 tension band.}
\pfield{Future Validation Path}{Final Fermilab result; J-PARC.}
\pfield{Falsification Criterion}{A final measured $a_\mu$ inconsistent with the spectral
value beyond the stated accuracy.}

\section{Oblique Parameters (S, T, U)}
\label{sec:oblique}

\predstart{AT-P017}{CONSISTENT}
\pfield{Title}{Oblique parameters from spectral structure.}
\pfield{Derived From}{Difference $\to$ Actualization $\to$ Spectrum (D96) $\to$ lightest-
octave fraction: $S = \mathrm{occ}_0/\Sigma m = 4/95 = 0.0421$ (dev.\ $5.3\%$);
$T = 2S = 8/95 = 0.0842$ (dev.\ $5.3\%$); $U = 0$ via the exact tree-level $\rho = 1$.}
\pfield{Prediction}{$S = 0.0421$, $T = 0.0842$ with $T = 2S$ exactly and $U = 0$
(claim classification: correspondence).}
\pfield{Experimental Test}{Electroweak global fit determination of $S,T,U$.}
\pfield{Current Evidence}{Consistent with the EW global fit beyond the masses and widths.}
\pfield{Future Validation Path}{Tightening of the EW fit at HL-LHC and future $e^+e^-$
colliders.}
\pfield{Falsification Criterion}{An EW global fit excluding $S = 0.0421$, $T = 0.0842$,
$U = 0$ at the stated accuracy.}

% ======================================================================
%  Part III — Gravity Predictions
% ======================================================================
\part{Gravity Predictions}

\section{Gravitational Sector}
\label{sec:gravity}

\predstart{AT-P018}{CONSISTENT}
\pfield{Title}{Newton constant from the D96 spectral content.}
\pfield{Derived From}{Difference $\to$ Actualization $\to$ Spectrum (D96) $\to$ occupation-
weighted content: $M_{\mathrm{Pl}} = v\cdot(\Sigma m\cdot\#g\cdot\mathrm{occ}_2)^3 =
1.22335\times10^{19}$ GeV (dev.\ $0.20\%$); $G = 1/M_{\mathrm{Pl}}^2 = 6.6476\times10^{-11}$
m$^3$ kg$^{-1}$ s$^{-2}$ (dev.\ $0.40\%$).}
\pfield{Prediction}{The Planck scale is the weak scale times the cube of the
occupation-weighted spectral content; the physical exponent is cubic (QG183)
(claim classification: calibration via the anchor $v$ plus SI conversion).}
\pfield{Experimental Test}{Measurement of $G$ (torsion balance) and $M_{\mathrm{Pl}}$.}
\pfield{Current Evidence}{Matches $G$ within $0.40\%$.}
\pfield{Future Validation Path}{Improved $G$ determinations (CODATA).}
\pfield{Falsification Criterion}{A measured $G$ inconsistent with the spectral construction
beyond the stated accuracy.}

\predstart{AT-P019}{CONSISTENT}
\pfield{Title}{Gravitational time dilation (GPS).}
\pfield{Derived From}{Difference $\to$ Actualization $\to$ counting measure $\rho$ $\to$
redshift law: clock rate $d\tau/dt = \rho^{1/d} = \sqrt{-g_{00}}$; GPS offset
$+38.5\ \mu$s/day vs.\ observed $+38.6\ \mu$s/day (dev.\ $0.2\%$).}
\pfield{Prediction}{Gravitational time dilation is the counting-measure redshift law; the
GPS rate offset is $-4.465\times10^{-10}$ (claim classification: correspondence).}
\pfield{Experimental Test}{GPS satellite clock comparisons.}
\pfield{Current Evidence}{The net computed offset $+38.5\ \mu$s/day matches the observed
$+38.6\ \mu$s/day within $0.2\%$.}
\pfield{Future Validation Path}{Continuing GPS and optical-clock tests.}
\pfield{Falsification Criterion}{A measured gravitational redshift inconsistent with the
counting-measure law beyond the stated accuracy.}

\predstart{AT-P020}{CONSISTENT}
\pfield{Title}{Frame dragging (Lense--Thirring).}
\pfield{Derived From}{Difference $\to$ tensor face $\psi$ against $\eta$ $\to$ restored
gravitomagnetic sector: $\Omega_{\mathrm{LT}} = G\bigl(3(\mathbf{J}\cdot\hat{\mathbf r})
\hat{\mathbf r}-\mathbf{J}\bigr)/(2c^2r^3)$; GP-B $41.1$ vs.\ $39.2$ mas/yr; LAGEOS $30.7$
vs.\ $\sim31$ mas/yr.}
\pfield{Prediction}{Frame dragging is a $\psi$-sector observable at full GR strength when
the tensor face is present (claim classification: correspondence).}
\pfield{Experimental Test}{GP-B and LAGEOS measurements of the Lense--Thirring precession.}
\pfield{Current Evidence}{GP-B $41.1$ vs.\ $39.2$ mas/yr ($+4.8\%$, inside the measured
$37.2\pm7.2$); LAGEOS $30.7$ vs.\ $\sim31$ ($-1.1\%$).}
\pfield{Future Validation Path}{Future gravity-mission measurements (e.g.\ proposed
LAGEOS-III-type and satellite tests).}
\pfield{Falsification Criterion}{A measured frame-dragging precession inconsistent with the
Lense--Thirring value at full GR strength.}

\section{Black-Hole Relations}
\label{sec:black-hole}

\predstart{AT-P021}{CONSISTENT}
\pfield{Title}{Mass-radius relation $M \propto R$.}
\pfield{Derived From}{Difference $\to$ Actualization $\to$ counting measure $\to$ per-octave
(logarithmic) deficit profile: $a \propto -1/r$, $GM_{\mathrm{eff}} \propto R$
(claim classification: correspondence; the flat-rotation-curve profile is a hosted input).}
\pfield{Prediction}{The observed $M \propto R$ relation follows from the counting measure:
the deficit is per-octave/log, giving the flat rotation-curve profile and
$GM_{\mathrm{eff}} \propto R$.}
\pfield{Experimental Test}{Galactic rotation curves and supermassive black-hole
mass--radius correlations.}
\pfield{Current Evidence}{Flat rotation curves are observed; $M \propto R$ is consistent
with the hosted profile.}
\pfield{Future Validation Path}{High-precision rotation curves (Rubin LSST) and horizon-scale
mass measurements (EHT).}
\pfield{Falsification Criterion}{A measured rotation-curve profile inconsistent with the
per-octave deficit profile across scales.}

\predstart{AT-P022}{CONSISTENT}
\pfield{Title}{Hawking temperature $T \propto 1/R$.}
\pfield{Derived From}{Same chain: with $S \propto R^{d-1}$ and $M \propto R$, the first law
gives $T \propto 1/R$; the $\psi$ sector contributes only a radius-independent prefactor.}
\pfield{Prediction}{The Hawking temperature scales as $T \propto 1/R$; the $T(R_1)/T(R_2)$
ratio is $\psi$-invariant (QG208).}
\pfield{Experimental Test}{Event-horizon-scale observations; astrophysical black-hole
cooling bounds.}
\pfield{Current Evidence}{Consistent with the standard $T = \hbar\kappa/(2\pi k_B)$
structure and current bounds.}
\pfield{Future Validation Path}{Gravitational-wave ringdown precision and horizon-scale
measurements.}
\pfield{Falsification Criterion}{A measured black-hole temperature inconsistent with
$T \propto 1/R$ beyond the stated accuracy.}

\section{Bekenstein Boundary}
\label{sec:bekenstein}

\predstart{AT-P023}{CONSISTENT (documented boundary)}
\pfield{Title}{Bekenstein structure $S \propto A$; exact $1/4$ coefficient is a boundary.}
\pfield{Derived From}{Difference $\to$ Actualization $\to$ counting measure: $S \propto A$
(QG12), $M \propto R$ (AT-P021), $T \propto 1/R$ (AT-P022) fully derived; the exact
coefficient $S = A/4$ is \emph{not} derivable from D96/TRM --- it requires the imported
quantum factor $T = \kappa/(2\pi)$ (QG185, QG196 impossibility proof).}
\pfield{Prediction}{The \emph{structure} of the Bekenstein law ($S \propto A$, $M \propto R$,
$T \propto 1/R$) is a derived consequence; the exact $1/4$ coefficient is a disclosed
boundary, not a derivation claim.}
\pfield{Experimental Test}{Tests of the area law and the entropy coefficient (astrophysical
and analogue systems).}
\pfield{Current Evidence}{The structure is consistent with black-hole thermodynamics; the
coefficient is not tested at the required precision.}
\pfield{Future Validation Path}{Quantum-gravity-sensitive observations; the coefficient
remains a documented boundary.}
\pfield{Falsification Criterion}{This entry is a boundary statement, not a falsifiable
prediction; it documents what the theory does not derive.}

% ======================================================================
%  Part IV — Cosmology Predictions
% ======================================================================
\part{Cosmology Predictions}

\section{$\Omega_\Lambda$}
\label{sec:omegal}

\predstart{AT-P024}{CONSISTENT}
\pfield{Title}{Vacuum density fraction.}
\pfield{Derived From}{Difference $\to$ Actualization $\to$ Spectrum (D96) $\to$ octave
occupancies $[4,4,87]$: $\Omega_\Lambda = I_{\mathrm{occ}}/\ln K_{\mathrm{oct}} = 0.6839$
(dev.\ $0.12\%$); $I_{\mathrm{occ}} = KL(p\Vert\mathrm{uniform}) = 0.7513$ nats;
$K_{\mathrm{oct}} = 3$ occupied octave classes.}
\pfield{Prediction}{The vacuum fraction is the information-density fraction of the D96
octave record (claim classification: correspondence).}
\pfield{Experimental Test}{Cosmological probes of the dark-energy fraction (CMB, supernovae,
BAO).}
\pfield{Current Evidence}{$\Omega_\Lambda = 0.6839$ matches the observed value within
$0.12\%$.}
\pfield{Future Validation Path}{DESI, Euclid, and Rubin LSST precision cosmology.}
\pfield{Falsification Criterion}{A measured $\Omega_\Lambda$ inconsistent with $0.6839$
beyond the stated accuracy.}

\section{$\Omega_m$}
\label{sec:omegam}

\predstart{AT-P025}{CONSISTENT}
\pfield{Title}{Matter density fraction.}
\pfield{Derived From}{Same chain: $\Omega_m = 1 - \Omega_\Lambda = 0.3161$ (dev.\ $0.26\%$).}
\pfield{Prediction}{The matter fraction is the complement of the octave-record information
fraction (claim classification: correspondence).}
\pfield{Experimental Test}{Cosmological probes of the matter density (CMB, BAO, galaxy
clustering, lensing).}
\pfield{Current Evidence}{$\Omega_m = 0.3161$ matches the observed value within $0.26\%$.}
\pfield{Future Validation Path}{DESI, Euclid, and Rubin LSST precision cosmology.}
\pfield{Falsification Criterion}{A measured $\Omega_m$ inconsistent with $0.3161$ beyond the
stated accuracy.}

\section{Acoustic Peaks}
\label{sec:peaks}

\predstart{AT-P026}{CONSISTENT}
\pfield{Title}{First acoustic peak existence (structural).}
\pfield{Derived From}{Difference $\to$ Actualization $\to$ Spectrum (D96) $\to$ recombination-
scale mode ladder: the fundamental doublet at $\omega_{\min} = 0.6216$ is isolated by the
dominant spectral gap (ln-gap $0.680$, largest of 94).}
\pfield{Prediction}{A first peak exists as a structural consequence: the fundamental doublet
is the maximally isolated low-frequency spectral feature and therefore the unique D96
candidate for a first peak (claim classification: theorem --- structural).}
\pfield{Experimental Test}{Observation of the CMB angular power spectrum.}
\pfield{Current Evidence}{The CMB first peak is observed at $\ell_1 \simeq 220$. Existence is
structural; the identification with the observed peak and its multipole position require
the $\ell$-axis normalization.}
\pfield{Future Validation Path}{CMB-S4 and Planck-precision acoustic spectra.}
\pfield{Falsification Criterion}{Absence of a first peak in the CMB angular spectrum.}

\predstart{AT-P027}{CONSISTENT}
\pfield{Title}{First peak location $\ell_1$ (fitted normalization).}
\pfield{Derived From}{Same chain: $\ell_1 = \Sigma m\cdot\ln(\mathrm{span})\cdot(5/4) =
220.48$ (dev.\ $0.008\%$).}
\pfield{Prediction}{$\ell_1 = 220.48$; the multiplier $5/4$ is a documented \emph{fit} for the
$\ell$-axis normalization (QG297), not part of the selection mechanism
(claim classification: fit).}
\pfield{Experimental Test}{Measurement of the first-peak multipole.}
\pfield{Current Evidence}{Observed $\ell_1 \simeq 220.48$; the value matches because the
$5/4$ factor is fitted to it.}
\pfield{Future Validation Path}{None beyond precision cosmology; the location is a documented
fit, not a derivation.}
\pfield{Falsification Criterion}{Not applicable --- the location is a documented fit, not a
falsifiable derivation.}

\predstart{AT-P028}{CONSISTENT}
\pfield{Title}{Peak ratios (spectral correspondences).}
\pfield{Derived From}{Same chain: $r_{21} = (\Sigma m-\#d)\cdot\mathrm{occ}_0/\mathrm{occ}_2 =
2.4368$ (dev.\ $0.035\%$), $r_{31} = \mathrm{span}/\sqrt{3} = 3.6965$ (dev.\ $0.058\%$).}
\pfield{Prediction}{The ratios of the second and third peak positions to the first are pure
spectral ratios, with no fitted constant (claim classification: correspondence).}
\pfield{Experimental Test}{Measurement of $\ell_2/\ell_1$ and $\ell_3/\ell_1$.}
\pfield{Current Evidence}{Matches observation within $\sim0.05\%$.}
\pfield{Future Validation Path}{CMB-S4 precision.}
\pfield{Falsification Criterion}{A measured peak ratio inconsistent with $2.4368$ or $3.6965$
beyond the stated accuracy.}

\predstart{AT-P029}{CONSISTENT}
\pfield{Title}{Scalar spectral index $n_s$ and running.}
\pfield{Derived From}{Same chain: $n_s = 0.96497$ (dev.\ $0.007\%$), running $= 0$
(QG237); seed spectrum is the Poisson counting variance, scale-free from critical
branching.}
\pfield{Prediction}{$n_s = 0.96497$ with zero running; the seed is the counting variance
(claim classification: correspondence).}
\pfield{Experimental Test}{Measurement of the scalar spectral index and its running from the
CMB power spectrum.}
\pfield{Current Evidence}{Planck measures $n_s = 0.9649 \pm 0.0042$; running consistent with
zero.}
\pfield{Future Validation Path}{CMB-S4.}
\pfield{Falsification Criterion}{A measured running significantly different from zero, or an
$n_s$ inconsistent with $0.96497$ beyond the stated accuracy.}

\section{Structure Formation}
\label{sec:structure}

\predstart{AT-P030}{CONSISTENT}
\pfield{Title}{Cosmological constant existence, sign, and scaling.}
\pfield{Derived From}{Difference $\to$ Actualization $\to$ critical branching vacuum $\to$
residual actualization pressure: at criticality ($\mu=1$) the Galton--Watson mean is
constant while the variance grows, so the realized vacuum never equals its uniform
expectation.}
\pfield{Prediction}{$\Lambda$ exists, is positive, and scales with the spectral/actualization
structure (claim classification: correspondence).}
\pfield{Experimental Test}{Measurement of the cosmological constant and its equation of
state.}
\pfield{Current Evidence}{$\Lambda > 0$ observed; $w \simeq -1$ consistent.}
\pfield{Future Validation Path}{DESI/Euclid dark-energy characterization.}
\pfield{Falsification Criterion}{A measured $\Lambda \le 0$, or a dark-energy dynamics
inconsistent with a constant residual pressure.}

\predstart{AT-P031}{PENDING}
\pfield{Title}{Structure formation from Q-event statistics.}
\pfield{Derived From}{Difference $\to$ Actualization $\to$ Q-event statistics and spectral
hierarchy $\to$ large-scale organization.}
\pfield{Prediction}{The large-scale organization of the universe follows the actualization
dynamics: structure formation is derived from the Q-event statistics and the spectral
hierarchy, with the initial conditions the uniform critical state and the information
content the departure that seeds the structure (claim classification: correspondence).}
\pfield{Experimental Test}{Galaxy-clustering statistics, the matter power spectrum, and the
CMB lensing signal.}
\pfield{Current Evidence}{Qualitative consistency with the standard structure-formation
picture; no discriminating quantitative test yet performed.}
\pfield{Future Validation Path}{DESI, Euclid, and Rubin LSST matter-power-spectrum
measurements.}
\pfield{Falsification Criterion}{A large-scale structure signal inconsistent with the
Q-event-statistics prediction at discriminating precision.}

\predstart{AT-P032}{CONSISTENT}
\pfield{Title}{Scale-free seed spectrum (Poisson counting variance).}
\pfield{Derived From}{Same chain: $\delta_i = 1/\sqrt{\langle N\rangle}$, scale-free
($n_s = 1$ at the seed level) from critical branching.}
\pfield{Prediction}{The seed power spectrum is the Poisson counting variance, scale-free from
critical branching (claim classification: correspondence).}
\pfield{Experimental Test}{Primordial power-spectrum shape from the CMB.}
\pfield{Current Evidence}{Consistent with the near-scale-invariant observed spectrum
($n_s = 0.96497$).}
\pfield{Future Validation Path}{CMB-S4 precision.}
\pfield{Falsification Criterion}{A primordial spectrum strongly inconsistent with the
counting-variance seed.}

% ======================================================================
%  Part V — Universality Predictions
% ======================================================================
\part{Universality Predictions}

\section{Operator Basis}
\label{sec:operators}

\predstart{AT-P033}{CONSISTENT}
\pfield{Title}{Four-operator basis universal across organized domains.}
\pfield{Derived From}{Difference $\to$ Actualization $\to$ Spectrum (D96) $\to$ operator
layer: Crowding (degeneracy), Compression (octaves), Beat (extent), Locking (gaps) as
spectrum projections.}
\pfield{Prediction}{The four-operator basis recurs as the structural read of organized
spectra in every organized domain (claim classification: correspondence).}
\pfield{Experimental Test}{Quantitative operator reads of organized systems: networks,
language, music, DNA, software, finance.}
\pfield{Current Evidence}{The four operators appear across the seven tested domains
(QG313); random spectra fail the basis while organized systems carry it (QG309, QG312).}
\pfield{Future Validation Path}{New organized-domain datasets.}
\pfield{Falsification Criterion}{An organized domain whose spectrum fails all four operators
at the quantitative level.}

\predstart{AT-P034}{CONSISTENT}
\pfield{Title}{No fifth operator.}
\pfield{Derived From}{Same chain: fifth-operator search (QG307) across unexplored domains.}
\pfield{Prediction}{No genuine fifth spectral operator exists beyond {Crowding, Compression,
Beat, Locking} with the universal MOMENT read-out (search-scoped, not an existence proof).}
\pfield{Experimental Test}{Search for an additional operator in organized-domain spectra.}
\pfield{Current Evidence}{No genuine fifth operator found in any searched domain; candidates
reduce to compositions or are not spectral reads.}
\pfield{Future Validation Path}{Continued domain expansion.}
\pfield{Falsification Criterion}{A genuine fifth operator (not a composition) in an
organized-domain spectrum.}

\section{Lock Law}
\label{sec:lock-law}

\predstart{AT-P035}{CONSISTENT}
\pfield{Title}{Lock law: universal structure, domain-specific values.}
\pfield{Derived From}{Difference $\to$ Actualization $\to$ Spectrum (D96) $\to$ Locking
operator $\to$ lock identities (normalized moment ratios).}
\pfield{Prediction}{Every organized domain carries stable, reproducible lock identities
(universal structure); the lock values are domain-specific fingerprints (claim
classification: correspondence).}
\pfield{Experimental Test}{Measurement of the moment/span, compression/count,
higher-moment, and $\sqrt{\text{moment}}$/span ratios across organized domains.}
\pfield{Current Evidence}{Stable lock structure in all seven tested domains (physics,
language, music, DNA, software, finance, networks); the D96 locks
$\mathrm{occMom}/\Sigma m = 20.0026$, $\Sigma m^2/\Sigma m = 2.4105$,
$\mathrm{occMom}/\Sigma m^2 = 8.2980$ follow the moment-chain identity exactly.}
\pfield{Future Validation Path}{New organized-domain cohorts.}
\pfield{Falsification Criterion}{An organized domain with no stable lock identity at the
quantitative level.}

\predstart{AT-P036}{CONSISTENT}
\pfield{Title}{Locks precede organization.}
\pfield{Derived From}{Same chain: early-lock prediction (QG315) on evolving systems.}
\pfield{Prediction}{In evolving systems, the lock structure appears before the full
organization develops, and systems with early lock structure reorganize less under
structural change (cohort-scoped on synthetic deterministic evolving-law cohorts).}
\pfield{Experimental Test}{Evolutionary cohorts (software history, wiki edits, citation
networks, language corpora).}
\pfield{Current Evidence}{Lock identity reaches half-strength earlier than maturity in most
systems; locked groups are 100\% small-reorganization.}
\pfield{Future Validation Path}{Real-world evolutionary datasets.}
\pfield{Falsification Criterion}{An evolving system whose lock structure lags the maturity
development.}

\section{Organization Transition}
\label{sec:transition}

\predstart{AT-P037}{CONSISTENT}
\pfield{Title}{Organization as a phase transition.}
\pfield{Derived From}{Same chain: organization-phase-transition audit (QG316) sweeping a
continuous organization parameter across the four regimes.}
\pfield{Prediction}{The binary operator basis completes discontinuously at the critical
parameter $g^\ast \simeq 0.31$, while the quantitative organization grows continuously;
the locks appear at the critical onset (cohort-scoped).}
\pfield{Experimental Test}{Controlled organization sweeps (white noise $\to$ strong
organization).}
\pfield{Current Evidence}{The all-four operator screen flips at $g^\ast \simeq 0.31$ and
persists for all stronger organization; the lock coherence emerges in the same window.}
\pfield{Future Validation Path}{Real-world organization transitions.}
\pfield{Falsification Criterion}{A continuous organization transition with no critical
threshold for the binary operator basis.}

\predstart{AT-P038}{CONSISTENT}
\pfield{Title}{Blind prediction of future organization class.}
\pfield{Derived From}{Same chain: blind-organization prediction (QG314) --- organization
score from lock structure only, rule fixed before the reveal.}
\pfield{Prediction}{The future organization class can be predicted from the early lock
structure with the prediction rule fixed in advance; the lock values predict the class, not
the strength within it.}
\pfield{Experimental Test}{Blind classification of future high-maturity systems from
early-stage lock coherence.}
\pfield{Current Evidence}{The early lock presence identifies the future top-class systems
exactly; the score separates the organized from the unorganized class.}
\pfield{Future Validation Path}{Prospective organizational datasets.}
\pfield{Falsification Criterion}{A future-class prediction failure beyond the stated
accuracy of the locked protocol.}

% ======================================================================
%  Part VI — Experimental Frontier
% ======================================================================
\part{Experimental Frontier}

\section{Predictions Awaiting Validation}
\label{sec:frontier}

The three pre-registered predictions P1, P2, and P3 are the experimental frontier of the
theory. They are frozen in the immutable registry (QG193); their values never change and
only their outcome states advance.

\predstart{AT-P039}{PENDING}
\pfield{Title}{P1 --- 106 GeV resonance.}
\pfield{Derived From}{Difference $\to$ Actualization $\to$ Spectrum (D96) $\to$ sector
ladder $\to$ missing-Z-anchor rung: $M_{106} = 7\cdot M_Z/6 = 7\cdot15.198$ GeV;
window $= M_{106} \pm$ (rung spacing)/2 $= 99$--$114$ GeV (stated
$98.79$--$113.99$).}
\pfield{Prediction}{A resonance at $106.39$ GeV central mass with window $99$--$114$ GeV;
production hierarchy 9 rungs $106.4\to263.4$ GeV; decay hierarchy unit quantum $15.20$ GeV
$\times 10$ + top $20.26$ GeV $\times 1$ (pre-registered, QG190; forbidden-input guard: no
ATLAS/CMS/fitted-mass constant).}
\pfield{Uncertainty}{$\pm7.60$ GeV (half the mean rung spacing); boson-anchor family agrees
within $0.74\%$.}
\pfield{Experimental Test}{LHC Run 3 and HL-LHC diphoton/diboson searches in the
$99$--$114$ GeV window.}
\pfield{Current Evidence}{PENDING: the window is neither confirmed nor excluded. The
$\sim95$ GeV scalar excess cluster aligns with the $91.19$ GeV rung (dev.\ $4.5\%$), not
with the $106.39$ GeV prediction; CMS/ATLAS full-Run-2 diphoton null searches cover
$106$ GeV with limits $15$--$102$ fb; LEP2 excludes SM-like $<114.4$ GeV but SM-strength
$hZZ$ only.}
\pfield{Future Validation Path}{HL-LHC $3000$ fb$^{-1}$ (projected $1$--$3$ fb sensitivity
in $100$--$106$ GeV); decisive.}
\pfield{Falsification Criterion}{No signal in statistically sensitive searches of the
$99$--$114$ GeV window (DISFAVORED/FALSIFIED).}

\predstart{AT-P040}{PENDING}
\pfield{Title}{P2 --- $0\nu\beta\beta$ effective Majorana mass.}
\pfield{Derived From}{Difference $\to$ Actualization $\to$ Spectrum (D96) $\to$ PMNS reads
(AT-P007) + absolute masses (AT-P009) + Majorana character (AT-P012):
$m_{\beta\beta} = \bigl|m_1c_{12}^2c_{13}^2 + m_2s_{12}^2c_{13}^2e^{i\alpha_2} +
m_3s_{13}^2e^{-2i\delta_\nu}\bigr| = 2.02$ meV (computed $2.0222$ meV), $\alpha_2 = \alpha_3
= 0$ (pre-registered, QG191; forbidden-input guard: no experimental limit or detector
sensitivity).}
\pfield{Uncertainty}{$\pm10\%$ ($1.8$--$2.2$ meV); dominated by $m_2 s_{12}^2 c_{13}^2 =
2.52$ meV, robust to the CP phase.}
\pfield{Experimental Test}{$0\nu\beta\beta$ search at nEXO / LEGEND-1000 (half-life
sensitivity $\sim10^{28}$ yr).}
\pfield{Current Evidence}{PENDING: below all current limits ($0.036$--$0.156$ eV).}
\pfield{Future Validation Path}{nEXO / LEGEND-1000 ton-scale reach.}
\pfield{Falsification Criterion}{A measured upper limit below $2.02$ meV (significant
exclusion) falsifies P2.}

\predstart{AT-P041}{PENDING}
\pfield{Title}{P3 --- Sector-ladder spectrum.}
\pfield{Derived From}{Difference $\to$ Actualization $\to$ Spectrum (D96) $\to$ ladder
radii $\times$ Z-anchor scale: $E_{\mathrm{rung}} = \mathrm{radius}\cdot(M_Z/6)$; unit
quantum $\Delta E = M_Z/6 = 15.20$ GeV, top quantum $20.26$ GeV (pre-registered, QG192;
forbidden-input guard: no collider bump or catalog).}
\pfield{Prediction}{Nine resonances: $106.39$ (primary) $\to 136.78 \to 151.98 \to 182.38
\to 197.58 \to 212.78 \to 227.97 \to 243.17 \to 263.43$ GeV, multiplicities unit
$\times 10$ + top $\times 1$, width scale $15.20$ GeV.}
\pfield{Uncertainty}{$\pm5\%$ per rung; boson-anchor family agrees within $0.74\%$.}
\pfield{Experimental Test}{High-energy searches at LHC Run 3 / HL-LHC and future
FCC-hh.}
\pfield{Current Evidence}{PENDING overall; the $151.98$ GeV rung matches the combined
$\sim152$ GeV diphoton excess (local $3.6\sigma$, global up to $5.4\sigma$; alignment
$0.0132\%$ deviation, look-elsewhere-corrected $z = 2.80\sigma$); SM anchors Z/H/t confirm
the ladder scale (QG200); eight rungs remain unprobed.}
\pfield{Future Validation Path}{HL-LHC confirmation of the $152$ GeV excess and full
Run-3 searches of the remaining rungs.}
\pfield{Falsification Criterion}{A sensitive search excludes any frozen rung (limit below
the rung energy) falsifies P3.}

\medskip
\noindent\hrule\par\smallskip
\noindent
\textbf{Registry rule:} no future phase may modify a registered prediction (QG193). Only the
outcome state may advance: PENDING $\to$ SUPPORTED $\to$ CONFIRMED, or PENDING $\to$
DISFAVORED $\to$ FALSIFIED. Frozen values never change.

% ======================================================================
%  Status Summary
% ======================================================================
\section*{Status Summary}
\label{sec:summary}

\begin{center}
\begin{tabular}{llp{4.6cm}p{4.2cm}}
\hline
\textbf{ID} & \textbf{Status} & \textbf{Prediction} & \textbf{Key value} \\
\hline
AT-P001 & CONSISTENT & Gauge structure $1+3+8$ & $12 = \deg C_{96}$ \\
AT-P002 & CONSISTENT & Gauge couplings & $137$; $3/95$; $8/64.08$; $0.2316$ \\
AT-P003 & CONSISTENT & Weak boson masses & $M_W$ 80.1, $M_Z$ 91.4 GeV \\
AT-P004 & CONSISTENT & Precision EW & $\sin^2\theta_{\mathrm{eff}}=0.23158$ \\
AT-P005 & CONSISTENT & CKM & dev.\ $0.58\%$ \\
AT-P006 & CONSISTENT & CKM CP & $\delta_{CP}=66.3^\circ$, $J=3.14\times10^{-5}$ \\
AT-P007 & CONSISTENT & PMNS & dev.\ $1.5\%$ \\
AT-P008 & CONSISTENT & Neutrino splittings & $7.607\times10^{-5}$; $2.44\times10^{-3}$ eV$^2$ \\
AT-P009 & PENDING & Absolute neutrino masses & $8.72$; $49.4$ meV; $\Sigma m_\nu=0.0581$ eV \\
AT-P010 & PENDING & Mass ordering (normal) & $m_1=0$ \\
AT-P011 & CONSISTENT & Quark masses & all six within $0.2\%$ \\
AT-P012 & PENDING & Majorana character & self-conjugate channel \\
AT-P013 & PENDING & $m_{\beta\beta}$ (P2) & $2.02$ meV \\
AT-P014 & CONSISTENT & Electron g-2 & $1.159655\times10^{-3}$ \\
AT-P015 & CONSISTENT & Electron anomaly-free & $\Delta a_e = 1.86\times10^{-13}$ \\
AT-P016 & CONSISTENT & Muon g-2 & $1.16644\times10^{-3}$ \\
AT-P017 & CONSISTENT & Oblique parameters & $S=0.0421$, $T=0.0842$, $U=0$ \\
AT-P018 & CONSISTENT & Newton constant & $6.6476\times10^{-11}$ m$^3$/kg/s$^2$ \\
AT-P019 & CONSISTENT & GPS time dilation & $+38.5\ \mu$s/day \\
AT-P020 & CONSISTENT & Frame dragging & GP-B $41.1$ mas/yr \\
AT-P021 & CONSISTENT & Mass--radius $M\propto R$ & $GM_{\mathrm{eff}}\propto R$ \\
AT-P022 & CONSISTENT & Hawking $T\propto 1/R$ & $\psi$-invariant ratio \\
AT-P023 & BOUNDARY & Bekenstein structure & $1/4$ coefficient not derived \\
AT-P024 & CONSISTENT & $\Omega_\Lambda$ & $0.6839$ \\
AT-P025 & CONSISTENT & $\Omega_m$ & $0.3161$ \\
AT-P026 & CONSISTENT & First peak existence & structural \\
AT-P027 & CONSISTENT & First peak location & $\ell_1 = 220.48$ (fit) \\
AT-P028 & CONSISTENT & Peak ratios & $2.4368$; $3.6965$ \\
AT-P029 & CONSISTENT & Spectral index & $n_s = 0.96497$, running $=0$ \\
AT-P030 & CONSISTENT & $\Lambda$ existence/sign & $\Lambda>0$ \\
AT-P031 & PENDING & Structure formation & Q-event statistics \\
AT-P032 & CONSISTENT & Seed spectrum & scale-free counting variance \\
AT-P033 & CONSISTENT & Four-operator universality & seven domains \\
AT-P034 & CONSISTENT & No fifth operator & search-scoped \\
AT-P035 & CONSISTENT & Lock law & universal structure, domain values \\
AT-P036 & CONSISTENT & Locks precede organization & cohort-scoped \\
AT-P037 & CONSISTENT & Organization transition & $g^\ast \simeq 0.31$ \\
AT-P038 & CONSISTENT & Blind class prediction & locked protocol \\
AT-P039 & PENDING & P1 --- 106 GeV & $106.39$ GeV, window $99$--$114$ \\
AT-P040 & PENDING & P2 --- $0\nu\beta\beta$ & $m_{\beta\beta}=2.02$ meV \\
AT-P041 & PENDING & P3 --- sector ladder & $9$ rungs $106.4\to263.4$ GeV \\
\hline
\end{tabular}
\end{center}

\medskip
\noindent
\fcolorbox{black}{gray!15}{%
  \begin{minipage}{0.92\linewidth}
    \centering
    \textbf{Registry totals}\\[0.4em]
    \begin{tabular}{cccc}
      \textbf{CONSISTENT} & \textbf{PENDING} & \textbf{BOUNDARY} & \textbf{FALSIFIED}\\[0.2em]
      32 & 8 & 1 & 0 \\
    \end{tabular}
  \end{minipage}}
\par
\smallskip
\noindent
No registered prediction is currently falsified. The three pre-registered frontier
predictions (AT-P039--AT-P041) are the experimental targets.

% ======================================================================
%  References
% ======================================================================
\section*{References}
\label{sec:references}

\begin{thebibliography}{99}
\bibitem{monograph} The Actualization Theory --- \emph{A Reconstruction of Physics from
Difference, Actualization and Spectrum}, Monograph V2.0, Docs/Publication/V2.0.
\bibitem{registry} AT-QG Phase 193, \emph{Prediction Registry Lock},
Docs/Research/ATQG\_PredictionRegistry.md and Docs/ATQG\_Predictions.md (immutable).
\bibitem{outcomes} AT-QG Phase 202, \emph{Prediction Outcome Dashboard},
Docs/ATQG\_PredictionOutcomes.md.
\bibitem{pre1} AT-QG Phase 190, \emph{Pre-Registered 106 GeV Resonance},
Docs/Research/ATQG\_PreRegistered106GeV.md.
\bibitem{pre2} AT-QG Phase 191, \emph{Pre-Registered $0\nu\beta\beta$},
Docs/Research/ATQG\_PreRegisteredMbb.md.
\bibitem{pre3} AT-QG Phase 192, \emph{Pre-Registered Sector-Ladder Spectrum},
Docs/Research/ATQG\_PreRegisteredLadderSpectrum.md.
\bibitem{class} AT-QG Claim Classification Registry, Docs/Research/ATQG\_ClaimClassificationRegistry.md.
\bibitem{auditmil} AT-QG Audit Milestone V2.0, Docs/Research/ATQG\_AUDIT\_MILESTONE\_V20.md.
\bibitem{release} AT-QG Release Check, Docs/Research/ATQG\_RELEASE\_CHECK.md.
\bibitem{qg188} AT-QG Phase 188, \emph{Prediction Audit}, Docs/Research.
\bibitem{qg188a} AT-QG Phase 188A, \emph{106 GeV Resonance Evidence Audit}, Docs/Research.
\bibitem{qg199} AT-QG Phase 199, \emph{P1 Evidence Update}, Docs/Research.
\bibitem{qg200} AT-QG Phase 200, \emph{Sector Ladder Evidence Audit}, Docs/Research.
\bibitem{qg201} AT-QG Phase 201, \emph{Ladder Statistics Audit}, Docs/Research.
\bibitem{qg203} AT-QG Phase 203, \emph{Absolute Neutrino Mass Origin}, Docs/Research.
\bibitem{valid001} AT-VALID001, \emph{Remaining Untested Items Audit}, Docs/Research/ATVALID\_UntestedItemsAudit.md.
\end{thebibliography}

\end{document}
